% Options for packages loaded elsewhere
\PassOptionsToPackage{unicode}{hyperref}
\PassOptionsToPackage{hyphens}{url}
%
\documentclass[
]{article}
\usepackage{amsmath,amssymb}
\usepackage{iftex}
\ifPDFTeX
  \usepackage[T1]{fontenc}
  \usepackage[utf8]{inputenc}
  \usepackage{textcomp} % provide euro and other symbols
\else % if luatex or xetex
  \usepackage{unicode-math} % this also loads fontspec
  \defaultfontfeatures{Scale=MatchLowercase}
  \defaultfontfeatures[\rmfamily]{Ligatures=TeX,Scale=1}
\fi
\usepackage{lmodern}
\ifPDFTeX\else
  % xetex/luatex font selection
\fi
% Use upquote if available, for straight quotes in verbatim environments
\IfFileExists{upquote.sty}{\usepackage{upquote}}{}
\IfFileExists{microtype.sty}{% use microtype if available
  \usepackage[]{microtype}
  \UseMicrotypeSet[protrusion]{basicmath} % disable protrusion for tt fonts
}{}
\makeatletter
\@ifundefined{KOMAClassName}{% if non-KOMA class
  \IfFileExists{parskip.sty}{%
    \usepackage{parskip}
  }{% else
    \setlength{\parindent}{0pt}
    \setlength{\parskip}{6pt plus 2pt minus 1pt}}
}{% if KOMA class
  \KOMAoptions{parskip=half}}
\makeatother
\usepackage{xcolor}
\setlength{\emergencystretch}{3em} % prevent overfull lines
\providecommand{\tightlist}{%
  \setlength{\itemsep}{0pt}\setlength{\parskip}{0pt}}
\setcounter{secnumdepth}{-\maxdimen} % remove section numbering
\ifLuaTeX
  \usepackage{selnolig}  % disable illegal ligatures
\fi
\IfFileExists{bookmark.sty}{\usepackage{bookmark}}{\usepackage{hyperref}}
\IfFileExists{xurl.sty}{\usepackage{xurl}}{} % add URL line breaks if available
\urlstyle{same}
\hypersetup{
  hidelinks,
  pdfcreator={LaTeX via pandoc}}

\author{}
\date{}

\begin{document}

\hypertarget{zero-shot-llm-translation-for-ancient-greek-lexicography-a-production-grade-human-in-the-loop-pipeline}{%
\section{Zero-Shot LLM Translation for Ancient Greek Lexicography: A
Production-Grade Human-in-the-Loop
Pipeline}\label{zero-shot-llm-translation-for-ancient-greek-lexicography-a-production-grade-human-in-the-loop-pipeline}}

\hypertarget{abstract}{%
\subsection{Abstract}\label{abstract}}

We present a production-oriented pipeline for generating and reviewing
English translations of ancient Greek lexicographic entries from
\emph{Stephanos of Byzantium}. The system uses off-the-shelf LLMs in
zero-shot mode (no fine-tuning) and couples them with strict OCR JSON
validation, database-backed provenance, prompt versioning, and human
review tooling. Beyond translation, the same architecture extracts cited
authors/works/fragments, people/deities/ethnic groups, aliases,
etymologies, and geospatial place links; it also generates
research-facing web pages and a PDF book. We focus on methodology rather
than model novelty: how to make LLM translation dependable enough for
large-scale scholarly workflows. We also introduce a structured
comparison layer between Billerbeck-based Greek and Meineke-linked
Greek, including translation-impact classification
(\texttt{likely\_different\_translation},
\texttt{probably\_same\_translation}, \texttt{uncertain}). Results from
operational use show that many observed textual differences are
mechanical and non-translation-changing, while a smaller subset likely
affects translation decisions. Finally, we outline a "Narrative
Learning" loop in which recurrent human-AI correction deltas are
distilled into explicit prompt updates. We release a reproducible
architecture pattern for digital humanities tasks where correctness is
cumulative and expert correction is mandatory.

\hypertarget{1-introduction}{%
\subsection{1. Introduction}\label{1-introduction}}

Large language models can generate fluent translations for historical
languages, but production use in scholarship fails when systems optimize
for output volume rather than correctness control. In ancient text
workflows, three constraints dominate:

\begin{itemize}
\tightlist
\item
  OCR noise propagates quickly.
\item
  Editorial traditions differ.
\item
  Human correction remains essential.
\end{itemize}

Our goal is to design a system that takes these constraints as
first-class requirements.

We report a deployed pipeline for the \emph{Ethnika} corpus that:

\begin{itemize}
\tightlist
\item
  ingests page-linked image data,
\item
  runs OCR with strict JSON output requirements,
\item
  assembles lemma-level records,
\item
  translates using zero-shot prompts,
\item
  supports human correction and review in a dedicated interface,
\item
  analyzes cross-edition textual differences for likely translation
  impact,
\item
  derives structured scholarly metadata (sources, works, fragments,
  entities, aliases, etymologies, places).
\end{itemize}

The key contribution is an end-to-end operational method that makes
zero-shot translation usable under expert supervision.

\hypertarget{2-task-and-data-model}{%
\subsection{2. Task and Data Model}\label{2-task-and-data-model}}

\hypertarget{21-unit-of-work}{%
\subsubsection{2.1 Unit of Work}\label{21-unit-of-work}}

The primary unit is a lemma record linked to one or more source page
images. Each lemma contains:

\begin{itemize}
\tightlist
\item
  machine OCR Greek text,
\item
  optional human-corrected Greek text,
\item
  machine translation outputs (versioned),
\item
  optional human translations,
\item
  review metadata.
\end{itemize}

\hypertarget{22-provenance-requirements}{%
\subsubsection{2.2 Provenance
Requirements}\label{22-provenance-requirements}}

Each stage writes explicit status flags and timestamps. Processing is
idempotent by default: only records with \texttt{processed=0} are acted
on. This minimizes accidental reprocessing and supports crash-safe
reruns.

\hypertarget{23-editorial-alignment}{%
\subsubsection{2.3 Editorial Alignment}\label{23-editorial-alignment}}

The system links assembled lemmas to Meineke headword paragraphs using
prioritized identifiers (\texttt{nodegoat\_id}, \texttt{billerbeck\_id},
\texttt{meineke\_id}). This enables downstream comparison even when base
working text comes from Billerbeck OCR. Each lemma also carries version
metadata (e.g., \texttt{epitome}, \texttt{parisinus},
\texttt{synthetic}), which is preserved through analysis and publication
layers.

\hypertarget{3-system-architecture}{%
\subsection{3. System Architecture}\label{3-system-architecture}}

\hypertarget{31-stage-a-ingestion}{%
\subsubsection{3.1 Stage A: Ingestion}\label{31-stage-a-ingestion}}

HTML/image registration and extraction populate image rows in
PostgreSQL. Inputs remain immutable; metadata and processing fields are
appended.

\hypertarget{32-stage-b-ocr}{%
\subsubsection{3.2 Stage B: OCR}\label{32-stage-b-ocr}}

OCR writes strict JSON to \texttt{images.lemma\_json}. Failures are
non-destructive:

\begin{itemize}
\tightlist
\item
  malformed JSON =\textgreater{} no \texttt{processed=1} update,
\item
  missing image =\textgreater{} hard error,
\item
  API failure =\textgreater{} bubble to batch runner.
\end{itemize}

\hypertarget{33-stage-c-assembly}{%
\subsubsection{3.3 Stage C: Assembly}\label{33-stage-c-assembly}}

OCR JSON is normalized into lemma-level structures. Human correction
fields are separate from machine fields.

\hypertarget{34-stage-d-translation}{%
\subsubsection{3.4 Stage D: Translation}\label{34-stage-d-translation}}

Translation uses zero-shot prompt templates with version tracking
(\texttt{translation\_prompt\_version}). New prompt versions trigger
targeted retranslation priority while preserving human edits.

\hypertarget{35-stage-e-review-and-publication}{%
\subsubsection{3.5 Stage E: Review and
Publication}\label{35-stage-e-review-and-publication}}

A CGI review interface displays source images, Billerbeck Greek views,
and cross-edition comparison status. Public/protected web pages are
generated from DB snapshots for transparency and operations monitoring.

\hypertarget{36-stage-f-structured-enrichment-modules}{%
\subsubsection{3.6 Stage F: Structured Enrichment
Modules}\label{36-stage-f-structured-enrichment-modules}}

After translation, additional forced-schema extraction jobs populate
reusable research tables:

\begin{itemize}
\tightlist
\item
  proper noun extraction with explicit \texttt{role} (\texttt{source} vs
  \texttt{entity}) and \texttt{noun\_type} (\texttt{person},
  \texttt{place}, \texttt{people}, \texttt{deity}, \texttt{other}),
\item
  source citation metadata (\texttt{citation}, \texttt{work\_title})
  enabling author/work/fragment indexing,
\item
  alias extraction (Stephanos-stated aliases plus rule-based spelling
  variants),
\item
  etymology extraction with controlled categories.
\end{itemize}

This design keeps all enrichment layers keyed to lemma IDs, so every
derived claim remains auditable against source Greek and OCR provenance.

\hypertarget{37-stage-g-geospatial-and-book-outputs}{%
\subsubsection{3.7 Stage G: Geospatial and Book
Outputs}\label{37-stage-g-geospatial-and-book-outputs}}

Place rows linked to coordinates and identifiers (Wikidata/Pleiades) are
rendered into an interactive Leaflet map. A separate publishing step
compiles a PDF book from database records, including indices for
sources, persons, places, peoples, and deities, plus map integration.

\hypertarget{4-llm-methodology}{%
\subsection{4. LLM Methodology}\label{4-llm-methodology}}

\hypertarget{41-zero-shot-constraint}{%
\subsubsection{4.1 Zero-Shot Constraint}\label{41-zero-shot-constraint}}

No model fine-tuning is used. We intentionally test a low-assumption
setup that many DH teams can replicate.

\hypertarget{42-prompt-governance}{%
\subsubsection{4.2 Prompt Governance}\label{42-prompt-governance}}

Prompts are stored as versioned rows. This creates:

\begin{itemize}
\tightlist
\item
  reproducibility across runs,
\item
  auditable changes in translation style/coverage,
\item
  controlled reprocessing policy.
\end{itemize}

\hypertarget{43-structured-difference-analysis}{%
\subsubsection{4.3 Structured Difference
Analysis}\label{43-structured-difference-analysis}}

For lemmas where normalized Billerbeck-vs-Meineke class is
\texttt{different}, we run forced-structure LLM analysis. Inputs are
pre-cleaned to remove non-semantic artifacts:

\begin{itemize}
\tightlist
\item
  Billerbeck leading entry numbers,
\item
  citation-heavy parenthetical wrappers.
\end{itemize}

Outputs include:

\begin{itemize}
\tightlist
\item
  difference level,
\item
  mechanical pattern tags,
\item
  word-pair substitutions,
\item
  translation-impact class.
\end{itemize}

\hypertarget{5-evaluation-strategy}{%
\subsection{5. Evaluation Strategy}\label{5-evaluation-strategy}}

Our evaluation is workflow-centric rather than benchmark-centric.

\hypertarget{51-what-we-measure}{%
\subsubsection{5.1 What We Measure}\label{51-what-we-measure}}

\begin{itemize}
\tightlist
\item
  processing success/failure rates,
\item
  number and category of Billerbeck-vs-Meineke differences,
\item
  translation-impact distribution,
\item
  human correction/review coverage,
\item
  extraction coverage for sources/entities/aliases/etymologies,
\item
  category-level and version-aware statistical signals.
\end{itemize}

\hypertarget{52-category-and-version-analytics}{%
\subsubsection{5.2 Category and Version
Analytics}\label{52-category-and-version-analytics}}

The statistics layer computes both descriptive and model-based analyses:

\begin{itemize}
\tightlist
\item
  word-count distributions by lemma type and initial letter,
\item
  ridge-regression models using proper noun features,
\item
  per-category models (authors, historical figures, places, ethnic
  groups, deities),
\item
  split analyses for Parisinus versus epitomised entries,
\item
  etymology category distributions by version.
\end{itemize}

This framework allows us to ask not only "what was translated?" but also
"what kinds of material were emphasized or compressed across textual
traditions?"

\hypertarget{53-why-this-matters}{%
\subsubsection{5.3 Why This Matters}\label{53-why-this-matters}}

For this task, an aggregate BLEU-style metric would obscure the real
risk profile. The key risk is not average lexical mismatch; it is the
subset of entries where model output and textual base choices can
mislead downstream scholarship.

\hypertarget{6-observed-outcomes-operational}{%
\subsection{6. Observed Outcomes
(Operational)}\label{6-observed-outcomes-operational}}

From current operational runs:

\begin{itemize}
\tightlist
\item
  the majority of detected cross-edition differences are mechanical,
\item
  a smaller subset is classified as likely translation-changing,
\item
  uncertainty remains concentrated in not-yet-analyzed rows.
\end{itemize}

This validates a triage model:

\begin{itemize}
\tightlist
\item
  analyze all "different" rows over time,
\item
  prioritize rows likely to change translation,
\item
  route flagged rows to human review first.
\end{itemize}

\hypertarget{7-error-analysis}{%
\subsection{7. Error Analysis}\label{7-error-analysis}}

Common error classes include:

\begin{itemize}
\tightlist
\item
  OCR artifacts that survive token-level cleaning,
\item
  citation-marker confusion,
\item
  morphological drift in dense lexicographic formulae,
\item
  over-smoothing in English rendering.
\end{itemize}

Mitigations that worked in practice:

\begin{itemize}
\tightlist
\item
  strict JSON validation gates,
\item
  non-destructive failure handling,
\item
  explicit comparison pre-cleaning,
\item
  forced-schema LLM outputs,
\item
  prompt version audits.
\end{itemize}

\hypertarget{8-human-in-the-loop-design}{%
\subsection{8. Human-in-the-Loop
Design}\label{8-human-in-the-loop-design}}

The review interface is not an afterthought; it is the core quality
mechanism. Effective elements include:

\begin{itemize}
\tightlist
\item
  immediate access to source images,
\item
  side-by-side Billerbeck text states (raw OCR and comparison text),
\item
  explicit Meineke-vs-Billerbeck status labels,
\item
  direct edit links from analysis pages.
\end{itemize}

This reduces context-switching and speeds expert correction.

\hypertarget{9-narrative-learning-loop}{%
\subsection{9. Narrative Learning
Loop}\label{9-narrative-learning-loop}}

We formalize prompt improvement as a data-driven narrative loop rather
than parameter fine-tuning.

\hypertarget{91-input-signal}{%
\subsubsection{9.1 Input Signal}\label{91-input-signal}}

Primary signal is structured disagreement between:

\begin{itemize}
\tightlist
\item
  AI translation output,
\item
  human-corrected or human-reviewed translation fields.
\end{itemize}

Each disagreement is linked to lemma metadata, source Greek, and review
provenance.

\hypertarget{92-pattern-mining}{%
\subsubsection{9.2 Pattern Mining}\label{92-pattern-mining}}

At scheduled intervals, we sample reviewed deltas and ask an LLM to:

\begin{itemize}
\tightlist
\item
  identify recurrent error motifs,
\item
  separate lexical/semantic issues from stylistic preference,
\item
  propose concise correction heuristics with examples.
\end{itemize}

\hypertarget{93-prompt-update-protocol}{%
\subsubsection{9.3 Prompt Update
Protocol}\label{93-prompt-update-protocol}}

Heuristics are converted into a new prompt version only after human
approval. We then:

\begin{enumerate}
\def\labelenumi{\arabic{enumi}.}
\tightlist
\item
  run targeted retranslation on a validation subset,
\item
  compare correction rates and reviewer effort against the prior
  version,
\item
  promote or rollback based on observed impact.
\end{enumerate}

This produces a text-level learning trajectory that is interpretable and
auditable. It aligns with your "Narrative Learning" framing (arXiv
preprint citation to be inserted), where system improvement is encoded
as explicit narrative constraints.

\hypertarget{10-operational-lessons}{%
\subsection{10. Operational Lessons}\label{10-operational-lessons}}

\begin{enumerate}
\def\labelenumi{\arabic{enumi}.}
\tightlist
\item
  Idempotency is non-negotiable for long-running pipelines.
\item
  "Parse first, mark processed later" prevents silent queue loss.
\item
  Prompt versioning is as important as model versioning.
\item
  Editorial comparison should be modeled as data, not narrative notes.
\item
  Small UI changes in review tools can have large impact on annotation
  throughput.
\end{enumerate}

\hypertarget{11-limitations-and-threats-to-validity}{%
\subsection{11. Limitations and Threats to
Validity}\label{11-limitations-and-threats-to-validity}}

\begin{itemize}
\tightlist
\item
  No gold-standard aligned translation benchmark is yet complete.
\item
  Human review depth may vary by entry.
\item
  Translation-impact labels are model-assisted and still need expert
  confirmation.
\item
  Results are reported from one corpus/domain and may not transfer
  unchanged.
\end{itemize}

\hypertarget{12-ethics-and-scholarly-safety}{%
\subsection{12. Ethics and Scholarly
Safety}\label{12-ethics-and-scholarly-safety}}

We treat model output as provisional. The system is designed to prevent
hidden automation authority:

\begin{itemize}
\tightlist
\item
  explicit review states,
\item
  tracked human edits,
\item
  no automatic overwrite of curated fields,
\item
  public distinction between machine and human outputs.
\end{itemize}

For classical scholarship, this separation is critical.

\hypertarget{13-conclusion}{%
\subsection{13. Conclusion}\label{13-conclusion}}

A zero-shot LLM setup can support ancient Greek translation workflows at
scale if embedded in a strict, provenance-rich, human-supervised
architecture. Our contribution is a deployable pattern for digital
humanities teams: combine conservative data engineering with targeted
model usage, and treat review as a first-class system component.

\hypertarget{reproducibility-notes}{%
\subsection{Reproducibility Notes}\label{reproducibility-notes}}

To replicate the method:

\begin{enumerate}
\def\labelenumi{\arabic{enumi}.}
\tightlist
\item
  Use idempotent stage flags for each processing phase.
\item
  Require strict machine-readable OCR payloads.
\item
  Separate machine and human correction columns.
\item
  Version translation prompts and retranslation policy.
\item
  Add structured cross-edition comparison with pre-cleaning.
\item
  Expose review-centric pages with direct record-level edit links.
\end{enumerate}

\hypertarget{placeholder-references}{%
\subsection{Placeholder References}\label{placeholder-references}}

\begin{itemize}
\tightlist
\item
  Prior work on machine translation for historical languages.
\item
  LLM reliability and structured output control.
\item
  Human-in-the-loop annotation systems.
\item
  Digital philology workflows for classical corpora.
\item
  {[}Your Name{]}, "Narrative Learning" (arXiv preprint; add full
  citation details).
\end{itemize}

\end{document}
